%%%%%%%%%%%%%%%%
%%% exod 28 %%%
%%%%%%%%%%%%%%%%

\Chapter{28}

\Verse{1}
{
	\hP{וְאַתָּה הַקְרֵב אֵלֶיךָ אֶת אַהֲרֹן אָחִיךָ וְאֶת בָּנָיו אִתּוֹ מִתּוֹךְ בְּנֵי יִשְׂרָאֵל לְכַהֲנוֹ לִי אַהֲרֹן נָדָב וַאֲבִיהוּא אֶלְעָזָר וְאִיתָמָר בְּנֵי אַהֲרֹן׃}
}
{
	\eP{
		“And you, bring Aaron, your brother, forward to you,
		and his sons with him, from among the children of Israel,
		for him to function as a priest for me: Aaron, Nadab and
		Abihu, Eleazar and Ithamar, Aaron’s sons.
	}
}
{
}

\Verse{2}
{
	\hP{וְעָשִׂיתָ בִגְדֵי קֹדֶשׁ לְאַהֲרֹן אָחִיךָ לְכָבוֹד וּלְתִפְאָרֶת׃}
}
{
	\eP{
		And you shall make holy clothes for Aaron, your brother, for glory
		and for beauty.
	}
}
{
%	\footnote{
%		Footnote 28:2. for glory and for beauty. Beauty inspires. Building
%		beautiful places for the practice of religion is a valuable thing.
%		Of course this does not mean building great edifices at the expense
%		of the starving masses, nor does it mean focusing on the outer
%		trappings and missing the content and spirit that they serve. There
%		must be balance—wisdom. But we must recognize the value of art and
%		beauty: the building, the priests’ clothing, the music, the smells,
%		the tastes. Religion is not the enemy of the senses.
%	}
}

\Verse{3}
{
	\hP{וְאַתָּה תְּדַבֵּר אֶל כּל חַכְמֵי לֵב אֲשֶׁר מִלֵּאתִיו רוּחַ חכְמָה וְעָשׂוּ אֶת בִּגְדֵי אַהֲרֹן לְקַדְּשׁוֹ לְכַהֲנוֹ לִי׃}
}
{
	\eP{
		And you shall speak to all the wise of heart whom I have
		filled with a spirit of wisdom, that they shall make
		Aaron’s clothes, to sanctify him for him to function as a
		priest for me.
	}
}
{
}

\Verse{4}
{
	\hP{וְאֵלֶּה הַבְּגָדִים אֲשֶׁר יַעֲשׂוּ חֹשֶׁן וְאֵפוֹד וּמְעִיל וּכְתֹנֶת תַּשְׁבֵּץ מִצְנֶפֶת וְאַבְנֵט וְעָשׂוּ בִגְדֵי קֹדֶשׁ לְאַהֲרֹן אָחִיךָ וּלְבָנָיו לְכַהֲנוֹ לִי׃}
}
{
	\eP{
		And these are the clothes that they shall make:
		breastplate and ephod and robe and patterned coat, a
		headdress and a sash. And they shall make holy clothes for
		Aaron, your brother, and for his sons, for him to function
		as a priest for me.
	}
}
{
}

\Verse{5}
{
	\hP{וְהֵם יִקְחוּ אֶת הַזָּהָב וְאֶת הַתְּכֵלֶת וְאֶת הָאַרְגָּמָן וְאֶת תּוֹלַעַת הַשָּׁנִי וְאֶת הַשֵּׁשׁ׃}
}
{
	\eP{
		And they shall take the gold and the blue and the purple
		and the scarlet and the linen.
	}
}
{
}

\Verse{6}
{
	\hP{וְעָשׂוּ אֶת הָאֵפֹד זָהָב תְּכֵלֶת וְאַרְגָּמָן תּוֹלַעַת שָׁנִי וְשֵׁשׁ משְׁזָר מַעֲשֵׂה חֹשֵׁב׃}
}
{
	\eP{
		“And they shall make the ephod: gold, blue, and purple,
		scarlet, and woven linen, designer’s work.
	}
}
{
}

\Verse{7}
{
	\hP{שְׁתֵּי כְתֵפֹת חֹבְרֹת יִהְיֶה לּוֹ אֶל שְׁנֵי קְצוֹתָיו וְחֻבָּר׃}
}
{
	\eP{
		It shall have two connected shoulder-pieces at its two
		ends, and it will be connected.
	}
}
{
}

\Verse{8}
{
	\hP{וְחֵשֶׁב אֲפֻדָּתוֹ אֲשֶׁר עָלָיו כְּמַעֲשֵׂהוּ מִמֶּנּוּ יִהְיֶה זָהָב תְּכֵלֶת וְאַרְגָּמָן וְתוֹלַעַת שָׁנִי וְשֵׁשׁ משְׁזָר׃}
}
{
	\eP{
		And the ephod’s designed belt that is on it shall be like
		its work and be a part of it: gold, blue, and purple, and
		scarlet and woven linen.
	}
}
{
}

\Verse{9}
{
	\hP{וְלָקַחְתָּ אֶת שְׁתֵּי אַבְנֵי שֹׁהַם וּפִתַּחְתָּ עֲלֵיהֶם שְׁמוֹת בְּנֵי יִשְׂרָאֵל׃}
}
{
	\eP{
		And you shall take two onyx stones and inscribe the names
		of the children of Israel on them:
	}
}
{
}

\Verse{10}
{
	\hP{שִׁשָּׁה מִשְּׁמֹתָם עַל הָאֶבֶן הָאֶחָת וְאֶת שְׁמוֹת הַשִּׁשָּׁה הַנּוֹתָרִים עַל הָאֶבֶן הַשֵּׁנִית כְּתוֹלְדֹתָם׃}
}
{
	\eP{
		six of their names on one stone and the names of the
		remaining six on the second stone, according to their
		birth orders.
	}
}
{
}

\Verse{11}
{
	\hP{מַעֲשֵׂה חָרַשׁ אֶבֶן פִּתּוּחֵי חֹתָם תְּפַתַּח אֶת שְׁתֵּי הָאֲבָנִים עַל שְׁמֹת בְּנֵי יִשְׂרָאֵל מֻסַבֹּת מִשְׁבְּצוֹת זָהָב תַּעֲשֶׂה אֹתָם׃}
}
{
	\eP{
		By a stone engraver’s work you shall inscribe the two
		stones with signet inscriptions according to the names of
		the children of Israel. You shall make them surrounded by
		settings of gold.
	}
}
{
}

\Verse{12}
{
	\hP{וְשַׂמְתָּ אֶת שְׁתֵּי הָאֲבָנִים עַל כִּתְפֹת הָאֵפֹד אַבְנֵי זִכָּרֹן לִבְנֵי יִשְׂרָאֵל וְנָשָׂא אַהֲרֹן אֶת שְׁמוֹתָם לִפְנֵי יְהֹוָה עַל שְׁתֵּי כְתֵפָיו לְזִכָּרֹן׃}
}
{
	\eP{
		And you shall set the two stones on the ephod’s
		shoulder-pieces, commemorative stones for the children of
		Israel, and Aaron shall carry their names in front of YHWH
		on his two shoulders as a commemoration.
	}
}
{
}

\Verse{13}
{
	\hP{וְעָשִׂיתָ מִשְׁבְּצֹת זָהָב׃}
}
{
	\eP{“And you shall make settings of gold}
}
{
}

\Verse{14}
{
	\hP{וּשְׁתֵּי שַׁרְשְׁרֹת זָהָב טָהוֹר מִגְבָּלֹת תַּעֲשֶׂה אֹתָם מַעֲשֵׂה עֲבֹת וְנָתַתָּה אֶת שַׁרְשְׁרֹת הָעֲבֹתֹת עַל הַמִּשְׁבְּצֹת׃}
}
{
	\eP{
		and two chains of pure gold. You shall make them twisted,
		rope work, and put the rope chains on the settings.
	}
}
{
}

\Verse{15}
{
	\hP{וְעָשִׂיתָ חֹשֶׁן מִשְׁפָּט מַעֲשֵׂה חֹשֵׁב כְּמַעֲשֵׂה אֵפֹד תַּעֲשֶׂנּוּ זָהָב תְּכֵלֶת וְאַרְגָּמָן וְתוֹלַעַת שָׁנִי וְשֵׁשׁ משְׁזָר תַּעֲשֶׂה אֹתוֹ׃}
}
{
	\eP{
		“And you shall make a breastplate of judgment, designer’s
		work. You shall make it like the work of the ephod. You
		shall make it gold, blue, and purple and scarlet and woven
		linen.
	}
}
{
}

\Verse{16}
{
	\hP{רָבוּעַ יִהְיֶה כָּפוּל זֶרֶת ארְכּוֹ וְזֶרֶת רחְבּוֹ׃}
}
{
	\eP{
		It shall be square when doubled, its length a span and its
		width a span.
	}
}
{
}

\Verse{17}
{
	\hP{וּמִלֵּאתָ בוֹ מִלֻּאַת אֶבֶן אַרְבָּעָה טוּרִים אָבֶן טוּר אֹדֶם פִּטְדָה וּבָרֶקֶת הַטּוּר הָאֶחָד׃}
}
{
	\eP{
		And you shall put a mounting of stones in it: four rows of stone. A
		row of a carnelian, a topaz, and an emerald: the one row.
	}
}
{
}

\Verse{18}
{
	\hP{וְהַטּוּר הַשֵּׁנִי נֹפֶךְ סַפִּיר וְיָהֲלֹם׃}
}
{
	\eP{And the second row: a ruby, a sapphire, and a diamond.}
}
{
}

\Verse{19}
{
	\hP{וְהַטּוּר הַשְּׁלִישִׁי לֶשֶׁם שְׁבוֹ וְאַחְלָמָה׃}
}
{
	\eP{And the third row: a jacinth, an agate, and an amethyst.}
}
{
}

\Verse{20}
{
	\hP{וְהַטּוּר הָרְבִיעִי תַּרְשִׁישׁ וְשֹׁהַם וְיָשְׁפֵה מְשֻׁבָּצִים זָהָב יִהְיוּ בְּמִלּוּאֹתָם׃}
}
{
	\eP{
		And the fourth row: a beryl and an onyx and a jasper. They
		shall be set in gold in their mountings.
	}
}
{
}

\Verse{21}
{
	\hP{וְהָאֲבָנִים תִּהְיֶיןָ עַל שְׁמֹת בְּנֵי יִשְׂרָאֵל שְׁתֵּים עֶשְׂרֵה עַל שְׁמֹתָם פִּתּוּחֵי חוֹתָם אִישׁ עַל שְׁמוֹ תִּהְיֶיןָ לִשְׁנֵי עָשָׂר שָׁבֶט׃}
}
{
	\eP{
		And the stones shall be with the names of the children of
		Israel, twelve with their names; they shall be signet
		engravings, each with its name, for twelve tribes.
	}
}
{
}

\Verse{22}
{
	\hP{וְעָשִׂיתָ עַל הַחֹשֶׁן שַׁרְשֹׁת גַּבְלֻת מַעֲשֵׂה עֲבֹת זָהָב טָהוֹר׃}
}
{
	\eP{
		And you shall make twisted chains, rope work, of pure
		gold, on the breastplate.
	}
}
{
}

\Verse{23}
{
	\hP{וְעָשִׂיתָ עַל הַחֹשֶׁן שְׁתֵּי טַבְּעוֹת זָהָב וְנָתַתָּ אֶת שְׁתֵּי הַטַּבָּעוֹת עַל שְׁנֵי קְצוֹת הַחֹשֶׁן׃}
}
{
	\eP{
		And you shall make two rings of gold on the breastplate
		and put the two rings on the two ends of the breastplate
	}
}
{
}

\Verse{24}
{
	\hP{וְנָתַתָּה אֶת שְׁתֵּי עֲבֹתֹת הַזָּהָב עַל שְׁתֵּי הַטַּבָּעֹת אֶל קְצוֹת הַחֹשֶׁן׃}
}
{
	\eP{
		and put the two ropes of gold on the two rings at the ends
		of the breastplate.
	}
}
{
}

\Verse{25}
{
	\hP{וְאֵת שְׁתֵּי קְצוֹת שְׁתֵּי הָעֲבֹתֹת תִּתֵּן עַל שְׁתֵּי הַמִּשְׁבְּצוֹת וְנָתַתָּה עַל כִּתְפוֹת הָאֵפֹד אֶל מוּל פָּנָיו׃}
}
{
	\eP{
		And you shall put the two ends of the two ropes on the two
		settings, and you shall put them on the ephod’s
		shoulder-pieces opposite the front of it.
	}
}
{
}

\Verse{26}
{
	\hP{וְעָשִׂיתָ שְׁתֵּי טַבְּעוֹת זָהָב וְשַׂמְתָּ אֹתָם עַל שְׁנֵי קְצוֹת הַחֹשֶׁן עַל שְׂפָתוֹ אֲשֶׁר אֶל עֵבֶר הָאֵפֹד בָּיְתָה׃}
}
{
	\eP{
		And you shall make two rings of gold and set them on the
		breastplate’s two ends on its side that is adjacent to the
		ephod, inward.
	}
}
{
}

\Verse{27}
{
	\hP{וְעָשִׂיתָ שְׁתֵּי טַבְּעוֹת זָהָב וְנָתַתָּה אֹתָם עַל שְׁתֵּי כִתְפוֹת הָאֵפוֹד מִלְּמַטָּה מִמּוּל פָּנָיו לְעֻמַּת מַחְבַּרְתּוֹ מִמַּעַל לְחֵשֶׁב הָאֵפוֹד׃}
}
{
	\eP{
		And you shall make two rings of gold and put them on the
		ephod’s two shoulder-pieces below, opposite the front of
		it, by its connection, above the ephod’s designed belt.
	}
}
{
}

\Verse{28}
{
	\hP{וְיִרְכְּסוּ אֶת הַחֹשֶׁן מִטַּבְּעֹתָו אֶל טַבְּעֹת הָאֵפוֹד בִּפְתִיל תְּכֵלֶת לִהְיוֹת עַל חֵשֶׁב הָאֵפוֹד וְלֹא יִזַּח הַחֹשֶׁן מֵעַל הָאֵפוֹד׃}
}
{
	\eP{
		And they shall attach the breastplate from its rings to
		the ephod’s rings with a blue string so it will be on the
		ephod’s designed belt and the breastplate will not be
		detached from the ephod.
	}
}
{
}

\Verse{29}
{
	\hP{וְנָשָׂא אַהֲרֹן אֶת שְׁמוֹת בְּנֵי יִשְׂרָאֵל בְּחֹשֶׁן הַמִּשְׁפָּט עַל לִבּוֹ בְּבֹאוֹ אֶל הַקֹּדֶשׁ לְזִכָּרֹן לִפְנֵי יְהֹוָה תָּמִיד׃}
}
{
	\eP{
		And Aaron shall carry the names of the children of Israel
		in the breastplate of judgment on his heart when he comes
		into the Holy, as a commemorative in front of YHWH always.
	}
}
{
}

\Verse{30}
{
	\hP{וְנָתַתָּ אֶל חֹשֶׁן הַמִּשְׁפָּט אֶת הָאוּרִים וְאֶת הַתֻּמִּים וְהָיוּ עַל לֵב אַהֲרֹן בְּבֹאוֹ לִפְנֵי יְהֹוָה וְנָשָׂא אַהֲרֹן אֶת מִשְׁפַּט בְּנֵי יִשְׂרָאֵל עַל לִבּוֹ לִפְנֵי יְהֹוָה תָּמִיד׃}
}
{
	\eP{
		And you shall put the Urim and Tummim in the breastplate
		of judgment, and they shall be on Aaron’s heart when he
		comes in front of YHWH, and Aaron shall carry the judgment
		of the children of Israel on his heart in front of YHWH
		always.
	}
}
{
%	\footnote{
%		Urim and Tummim. We have pondered for centuries what these are. They
%		have something to do with inquiring of God for an answer to a question.
%		They may contain letters that can spell out long answers, or they may
%		provide only a “yes” or “no,” or they may also provide a third option of
%		“no response.” They are mentioned four times in the Torah, twice in 1
%		Samuel (14:41, where the Greek text differs from the Masoretic Text by a
%		full line; and 25:6), and then never again in Israel’s history until the
%		narrative of Ezra and Nehemiah, where it is noted that answers through
%		the Urim and Tummim are not available (Ezra 2:63; Neh 7:65). We must
%		admit it: we just do not know what they are. What is important is that
%		several different biblical sources indicate that there was a belief that
%		it was possible to ask questions of God and get an answer, and that this
%		was done through a priest, not a prophet. It was a mechanism other than
%		prophecy to learn the will of God. And it is known as a practice only in
%		Israel’s earliest era.
%	}
}

\Verse{31}
{
	\hP{וְעָשִׂיתָ אֶת מְעִיל הָאֵפוֹד כְּלִיל תְּכֵלֶת׃}
}
{
	\eP{“And you shall make the robe of the ephod all of blue.}
}
{
}

\Verse{32}
{
	\hP{וְהָיָה פִי רֹאשׁוֹ בְּתוֹכוֹ שָׂפָה יִהְיֶה לְפִיו סָבִיב מַעֲשֵׂה אֹרֵג כְּפִי תַחְרָא יִהְיֶה לּוֹ לֹא יִקָּרֵעַ׃}
}
{
	\eP{
		And its head-opening shall be within it. It shall have a
		binding for its opening all around, weaver’s work. It
		shall be like the opening of a coat of mail for him; it
		will not be torn.
	}
}
{
}

\Verse{33}
{
	\hP{וְעָשִׂיתָ עַל שׁוּלָיו רִמֹּנֵי תְּכֵלֶת וְאַרְגָּמָן וְתוֹלַעַת שָׁנִי עַל שׁוּלָיו סָבִיב וּפַעֲמֹנֵי זָהָב בְּתוֹכָם סָבִיב׃}
}
{
	\eP{
		And you shall make on its skirts pomegranates of blue and
		purple and scarlet, on its skirts all around, and bells of
		gold within them all around:
	}
}
{
}

\Verse{34}
{
	\hP{פַּעֲמֹן זָהָב וְרִמּוֹן פַּעֲמֹן זָהָב וְרִמּוֹן עַל שׁוּלֵי הַמְּעִיל סָבִיב׃}
}
{
	\eP{
		a golden bell and a pomegranate, a golden bell and a
		pomegranate, on the robe’s skirts all around.
	}
}
{
}

\Verse{35}
{
	\hP{וְהָיָה עַל אַהֲרֹן לְשָׁרֵת וְנִשְׁמַע קוֹלוֹ בְּבֹאוֹ אֶל הַקֹּדֶשׁ לִפְנֵי יְהֹוָה וּבְצֵאתוֹ וְלֹא יָמוּת׃}
}
{
	\eP{
		And it shall be on Aaron for ministering: and the sound of
		it will be heard when he comes to the Holy in front of
		YHWH and when he goes out, so he will not die.
	}
}
{
}

\Verse{36}
{
	\hP{וְעָשִׂיתָ צִּיץ זָהָב טָהוֹר וּפִתַּחְתָּ עָלָיו פִּתּוּחֵי חֹתָם קֹדֶשׁ לַיהֹוָה׃}
}
{
	\eP{
		“And you shall make a plate of pure gold and make signet
		inscriptions on it: ‘Holiness to YHWH.’
	}
}
{
}

\Verse{37}
{
	\hP{וְשַׂמְתָּ אֹתוֹ עַל פְּתִיל תְּכֵלֶת וְהָיָה עַל הַמִּצְנָפֶת אֶל מוּל פְּנֵי הַמִּצְנֶפֶת יִהְיֶה׃}
}
{
	\eP{
		And you shall set it on a blue string, and it shall be on
		the headdress. It shall be opposite the front of the
		headdress.
	}
}
{
}

\Verse{38}
{
	\hP{וְהָיָה עַל מֵצַח אַהֲרֹן וְנָשָׂא אַהֲרֹן אֶת עֲוֺן הַקֳּדָשִׁים אֲשֶׁר יַקְדִּישׁוּ בְּנֵי יִשְׂרָאֵל לְכל מַתְּנֹת קדְשֵׁיהֶם וְהָיָה עַל מִצְחוֹ תָּמִיד לְרָצוֹן לָהֶם לִפְנֵי יְהֹוָה׃}
}
{
	\eP{
		And it shall be on Aaron’s forehead, and Aaron shall bear the crime
		of the holy things that the children of Israel will sanctify, for
		all their holy gifts. And it shall be on his forehead always for
		acceptance for them in front of YHWH.
	}
}
{
%	\footnote{
%		28:38,43. bear a crime. This difficult expression is thought to have
%		a range of meanings. In these passages it conveys that the priests
%		bear an enormous responsibility on account of their closeness to the
%		holy. Coming close to the holy objects or to the altar, they risk
%		violation of the sacred if they act incorrectly. The seriousness of
%		this matter will become manifest when two of Aaron’s sons violate
%		the holy and die (Leviticus 10). After the destruction of the first
%		Temple in Jerusalem and the loss of the ark, tablets, Tabernacle,
%		and Urim and Tummim, much of the feeling for the sacred must have
%		been lost. And since the destruction of the second Temple, our sense
%		of the sacred has diminished even more. So it has become ever harder
%		to comprehend this aspect of the biblical world and to appreciate
%		the power and awe that were associated with sacred objects, sacred
%		places, and the priesthood.
%	}
}

\Verse{39}
{
	\hP{וְשִׁבַּצְתָּ הַכְּתֹנֶת שֵׁשׁ וְעָשִׂיתָ מִצְנֶפֶת שֵׁשׁ וְאַבְנֵט תַּעֲשֶׂה מַעֲשֵׂה רֹקֵם׃}
}
{
	\eP{
		And you shall weave the coat in a pattern of linen and
		make a headdress of linen. And you shall make a sash,
		embroiderer’s work,
	}
}
{
}

\Verse{40}
{
	\hP{וְלִבְנֵי אַהֲרֹן תַּעֲשֶׂה כֻתֳּנֹת וְעָשִׂיתָ לָהֶם אַבְנֵטִים וּמִגְבָּעוֹת תַּעֲשֶׂה לָהֶם לְכָבוֹד וּלְתִפְאָרֶת׃}
}
{
	\eP{
		and you shall make coats for Aaron’s sons, and you shall
		make sashes for them, and you shall make hats for them—for
		glory and for beauty.
	}
}
{
}

\Verse{41}
{
	\hP{וְהִלְבַּשְׁתָּ אֹתָם אֶת אַהֲרֹן אָחִיךָ וְאֶת בָּנָיו אִתּוֹ וּמָשַׁחְתָּ אֹתָם וּמִלֵּאתָ אֶת יָדָם וְקִדַּשְׁתָּ אֹתָם וְכִהֲנוּ לִי׃}
}
{
	\eP{
		And you shall dress them—Aaron, your brother, and his sons
		with him—and you shall anoint them and fill their hand and
		sanctify them, and they shall function as priests for me.
	}
}
{
%	\footnote{
%		anoint them. The Hebrew word for anointing is mš. Its English form is
%		messiah. Later it is used for kings. Kings Saul and David are each
%		referred to as the messiah (1 Sam 12:3; 24:7; 2 Sam 19:22). And, later
%		still, Jewish ideas of a coming messiah (Hebrew mšîa) developed out of
%		this royal role of anointing. But here in its first occurrence in the
%		Torah, it is understood that the original use of anointing is for
%		ordaining the high priests. 28:41; 29:9. fill their hand. In the Torah
%		and elsewhere in the Bible, this expression, “to fill one’s hand,” means
%		to take on the role of priest (Lev 8:33; 16:32; 21:10; Num 3:3; 1 Kings
%		13:33). In Leviticus it is usually related to donning the priestly
%		garments.
%	}
}

\Verse{42}
{
	\hP{וַעֲשֵׂה לָהֶם מִכְנְסֵי בָד לְכַסּוֹת בְּשַׂר עֶרְוָה מִמּתְנַיִם וְעַד יְרֵכַיִם יִהְיוּ׃}
}
{
	\eP{
		And make shorts of linen for them, to cover naked flesh;
		they shall be from hips to thighs.
	}
}
{
}

\Verse{43}
{
	\hP{וְהָיוּ עַל אַהֲרֹן וְעַל בָּנָיו בְּבֹאָם אֶל אֹהֶל מוֹעֵד אוֹ בְגִשְׁתָּם אֶל הַמִּזְבֵּחַ לְשָׁרֵת בַּקֹּדֶשׁ וְלֹא יִשְׂאוּ עָוֺן וָמֵתוּ חֻקַּת עוֹלָם לוֹ וּלְזַרְעוֹ אַחֲרָיו׃}
}
{
	\eP{
		And they shall be on Aaron and on his sons when they come
		to the Tent of Meeting or when they come close to the
		altar to minister in the Holy, so they will not bear a
		crime and die. It is an eternal law for him and for his
		seed after him.
	}
}
{
}
